\section{Introduction}
\label{sec:intro}

Large language models are no longer used only as static generators that answer a single user query. They now drive autonomous agents that act in external environments over many steps: calling tools to retrieve information, operating a browser to complete a task, or coordinating with other agents on a larger workflow. An LLM-based agent runs in a loop of perception, decision, execution, and observation, and in doing so it continually produces interaction records---reasoning chains, tool calls, environment observations, and human feedback. These records are the raw material from which an agent improves its own decisions. We formalize each step as a tuple $e = (c, a, o, f)$, whose elements are the decision context, the agent's output, the environment's response, and an optional evaluative signal. The community has built a range of mechanisms that turn this raw material into reusable form: memory summarization distills raw trajectories into natural-language reflections; skill induction turns successful trajectories into reusable code; soft-prompt training compresses experience into continuous latent representations; reward modeling internalizes human preference annotations as a parametric evaluator; trajectory fine-tuning writes filtered successful trajectories into policy weights; and RLHF internalizes evaluative signals as decision preferences in those weights.

The questions that arise when building an agent system cut across these mechanisms. Given one successful trajectory, should it become a retrievable reflection note or a mountable soft prompt? To give an agent a particular skill, should that skill be written into the weights by trajectory fine-tuning, or kept as an externally auditable memory entry? The answers are scattered across the memory and post-training literatures, and the taxonomies in current use---organized by \emph{component} (memory, planning, tool use) or by \emph{technique} (RAG, SFT, RLHF)---do not address them directly.

These mechanisms are instances of one operation. Each re-represents experience and, in doing so, re-balances a common set of trade-offs across interpretability, inference cost, and editability. Consider a single successful trajectory. The most direct use is to place it in the prompt as a few-shot example---direct reuse, at the cost of occupying the context window on every call. Memory summarization distills the trajectory into a more compact reflection or rule that remains readable, editable natural language but is far denser, so it consumes much less context. Trajectory fine-tuning instead trains the trajectory into the model's weights: nothing is occupied in context at inference time, but the experience can no longer be inspected or revised entry by entry. Memory summarization, skill induction, reward modeling, trajectory fine-tuning, and self-generated training can all be read as instances of this operation between different source and target carriers. We call it \emph{Experience Transformation}.

Different transformation paths run between different carriers. We classify carriers by where experience exists within the model architecture, giving three top-level classes:
\begin{itemize}
  \item \textbf{Tokenized carriers}: experience that enters the forward pass explicitly, as a discrete token sequence. By degree of formalization, these divide into \emph{Narrative} forms (reflections, summaries, and other natural language) and \emph{Schematic} forms (code, workflows, knowledge graphs, and other structured artifacts).
  \item \textbf{Latent carriers}: intermediate representations that exist as continuous vectors or hidden states and participate directly in attention or hidden-state computation.
  \item \textbf{Parametric carriers}: experience folded implicitly into network weights. By functional role, these divide into \emph{Policy} weights---the actor that generates actions---and \emph{Evaluator} weights---the judge that assesses trajectories.
\end{itemize}

Migration among these three classes is the through-line of this survey. We organize the mechanisms scattered across memory, skill induction, alignment, and self-training into seven basic transformation pathways:
\begin{itemize}
  \item \textbf{Narrative Abstraction} (Narrative $\rightarrow$ Narrative): distilling raw trajectories into more compact natural-language forms---reflections, rules, summaries---within the same carrier.
  \item \textbf{Schematic Formalization} (Narrative $\rightarrow$ Schematic): converting natural-language experience into code, workflows, or knowledge graphs that can be executed symbolically or traversed.
  \item \textbf{Latent-Space Transformation} (Tokenized $\rightarrow$ Latent): converting discrete tokenized experience into continuous vectors that join the forward computation directly.
  \item \textbf{Evaluator Internalization} (Tokenized $\rightarrow$ Evaluator): folding trajectory experience into a parametric evaluation capability, such as a reward model, process reward model, or verifier.
  \item \textbf{Policy Internalization} (Tokenized $\rightarrow$ Policy): folding trajectory experience into the decision capability of policy weights.
  \item \textbf{Evaluator-Driven Optimization} (Evaluator $\rightarrow$ Policy): absorbing evaluator signals into the generative preferences of policy weights.
  \item \textbf{Parametric Externalization} (Parametric $\rightarrow$ Tokenized): externalizing implicit weight experience as new tokenized trajectories.
\end{itemize}

Beyond these single pathways, we identify \emph{composite pipelines} that chain several basic pathways together, and we treat them as a separate object of analysis.

This survey fills an analytical dimension that existing surveys leave uncovered. Surveys of agent architecture organize systems by component---memory, planning, tool use---but one kind of experience-transformation operation crosses those component boundaries: trajectory fine-tuning, for instance, belongs to neither memory nor planning and has no home in a component view. Surveys of memory systems characterize storage, retrieval, and forgetting within Tokenized carriers in depth, but do not treat experience that migrates across carriers. Alignment surveys follow the technical progression from RLHF to DPO and its variants along the single Evaluator $\rightarrow$ Policy pathway, and surveys of process reward models concentrate on the design and training of the Evaluator carrier itself. Each provides depth within one system architecture, one carrier, or one technique family; none takes the transformation of experience between carriers as a dimension of analysis in its own right. We build a unified analytical framework around experience transformation. Our contributions are:
\begin{itemize}
  \item \emph{(Contribution 1 --- to be finalized.)}
  \item \emph{(Contribution 2 --- to be finalized.)}
  \item \emph{(Contribution 3 --- to be finalized.)}
\end{itemize}

The rest of the survey is organized as follows. Section~\ref{sec:prelim} establishes the conceptual framework. Sections~\ref{sec:tok2tok}--\ref{sec:param-source} analyze the literature for the seven single pathways, grouped by source-carrier type. Section~\ref{sec:composite} examines composite pipelines, Section~\ref{sec:synthesis} carries out the cross-pathway synthesis and utilization analysis, and Section~\ref{sec:open} discusses open problems and future directions.
